\subsection{Sufficiency, Exponential Families, and Local Asymptotics}

The shape of a model determines which summaries of the data preserve information and which approximations become valid in large samples.

\begin{definition} \label{def:statistics-sufficiency-ancillarity} \hlt{Statistics, Sufficiency, and Ancillarity}\\
A \hlt{statistic} is a measurable map $T:\mathcal{X}\to\mathcal{T}$.  It is \hlt{sufficient} for $\theta$ if the conditional distribution of $X$ given $T(X)$ does not depend on $\theta$.  It is \hlt{ancillary} if the distribution of $T(X)$ is the same for every $\theta\in\Theta$.
\end{definition}

\begin{theorem} \label{thm:factorization-criterion-sufficiency} \hlt{Factorization Criterion}\\
Let $\{P_\theta:\theta\in\Theta\}$ be dominated by $\mu$ with densities $p_\theta$.  A statistic $T$ is sufficient if and only if there exist nonnegative measurable functions $g_\theta$ and $h$ such that
\[
p_\theta(x)=g_\theta(T(x))h(x)
\]
for every $\theta$ and for $\mu$-almost every $x$.
\begin{proof}
If the factorization holds, all parameter dependence in the density is carried by $T(x)$, so the conditional law inside a fiber of $T$ is proportional to $h(x)$ and is parameter-free.  Conversely, if the conditional law given $T$ is parameter-free, disintegrating $P_\theta$ into the marginal law of $T$ and the common conditional law gives the displayed factorization after choosing densities with respect to a dominating measure.
\end{proof}
\end{theorem}

\begin{definition} \label{def:exponential-family-natural-statistic} \hlt{Exponential Family}\\
A dominated model is a \hlt{$k$-parameter exponential family} if its densities can be written
\[
p_\eta(x)=\exp\!\left(\eta^\top T(x)-A(\eta)\right)h(x),
\qquad \eta\in\mathcal{N}\subseteq\R^k.
\]
Here $\eta$ is the \hlt{natural parameter}, $T(X)$ is the \hlt{natural statistic}, and $A(\eta)$ is the \hlt{log-partition function}.
\end{definition}

\begin{proposition} \label{prop:exponential-family-moments} \hlt{Moment Identities in Full Exponential Families}\\
Assume the natural parameter space is open and differentiation under the integral is valid.  Then
\[
\nabla A(\eta)=\mathbb{E}_\eta[T(X)],
\qquad
\nabla^2 A(\eta)=\operatorname{Cov}_\eta(T(X)).
\]
\begin{proof}
Since $1=\int \exp(\eta^\top T(x)-A(\eta))h(x)\,d\mu(x)$, differentiating once gives $\nabla A(\eta)=\mathbb{E}_\eta[T(X)]$.  Differentiating again gives the covariance matrix of $T(X)$.
\end{proof}
\end{proposition}

\begin{definition} \label{def:local-asymptotic-normality} \hlt{Local Asymptotic Normality}\\
A sequence of experiments $\{P_{\theta}^{(n)}:\theta\in\Theta\}$ is \hlt{locally asymptotically normal} at $\theta_0$ with information matrix $I(\theta_0)$ if, for fixed $h\in\R^d$,
\[
\log\frac{dP_{\theta_0+h/\sqrt{n}}^{(n)}}{dP_{\theta_0}^{(n)}}
=h^\top\Delta_n-\frac{1}{2}h^\top I(\theta_0)h+o_{P_{\theta_0}}(1),
\]
where $\Delta_n\Rightarrow\mathcal{N}(0,I(\theta_0))$ under $P_{\theta_0}^{(n)}$.
\end{definition}

\begin{remark} \label{rem:lan-bridge-to-estimation-tests} \hlt{Why Local Asymptotics Matter}\\
Local asymptotic normality is the bridge from finite-sample likelihoods to the normal approximations used by estimators, tests, and confidence regions.  It packages the central-limit behavior of the score into a model-level approximation.
\end{remark}
